\chapter{总结与展望}
\label{chap:7}
\chapterguide{总结只采用正文和实验能够支持的结论。贡献与不足须对应材料和实际分工，不使用未经比较的“首次”“全面无损”等表述。}

\section{研究工作总结}
\label{sec:7-1}

\subsection{差异分析与统一模型设计总结}
\label{subsec:7-1-1}
\outlinehint{概括参数含义、引用及截面组织的统一表达，指出五类对象之间的共用结构与作用。}

\subsection{转换方法与系统实现总结}
\label{subsec:7-1-2}
\outlinehint{总结坐标和方向处理、引用恢复、材料范围及扫掠联合重建，区分个人工作和协作集成。}

\subsection{实验结论与适用范围总结}
\label{subsec:7-1-3}
\outlinehint{依据第六章真实数据概括可编辑和几何保留程度，明确软件版本、路线、子类型和失败条件；当前留待实验完成。}

\section{研究不足与后续展望}
\label{sec:7-2}

\subsection{草图尺寸、外部关联与终止条件完善}
\label{subsec:7-2-1}
\outlinehint{围绕驱动尺寸、投影关联、复杂参考与不等量薄壁等现有不足提出可执行的改进方向。}

\subsection{复杂扫掠与统一验证体系扩展}
\label{subsec:7-2-2}
\outlinehint{讨论更复杂路径、截面控制和样条的后续研究，以及目标重读、完整参数比较和可复现实验集的完善。}
